% =====================================================================
%  preamble.tex — пакеты, стили и макросы учебного курса
%  Зона D (research/paper/**). Движок: pdflatex (T2A + babel), MiKTeX.
%  Выбор pdflatex, а не XeLaTeX: так же собран seed-документ проекта
%  (research/paper/problem_statement.tex), меньше внешних зависимостей.
% =====================================================================

% --- Кодировка и язык ---
\usepackage[T2A]{fontenc}
\usepackage[utf8]{inputenc}
\usepackage[english,russian]{babel}

% --- Математика ---
\usepackage{amsmath,amssymb,amsthm}
\usepackage{bm}

% --- Вёрстка ---
\usepackage[margin=24mm,bottom=28mm]{geometry}
\usepackage{microtype}
\usepackage{enumitem}
\usepackage{booktabs}
\usepackage{array}
\usepackage{graphicx}
\usepackage{xcolor}
\usepackage{caption}
\usepackage[most]{tcolorbox}
\usepackage{fancyhdr}
\usepackage{titlesec}
\usepackage{needspace}   % удержать врезку целиком: \needspace{N\baselineskip}

% --- Ссылки (последним, hyperref любит быть в конце) ---
\usepackage[colorlinks=true,linkcolor=cAccent,urlcolor=cAccent,citecolor=cAccent]{hyperref}

% =====================================================================
%  Палитра. Совпадает со стилем графиков figures/pstyle.py, чтобы текст
%  и иллюстрации читались как один документ.
% =====================================================================
\definecolor{cInk}{HTML}{0F172A}      % основной текст
\definecolor{cInk2}{HTML}{475569}     % вторичный текст
\definecolor{cAccent}{HTML}{0B5AA8}   % заголовки, ссылки
\definecolor{cS1}{HTML}{2A78D6}       % слот 1 графиков
\definecolor{cS2}{HTML}{EB6834}       % слот 2 графиков
\definecolor{cS3}{HTML}{1BAF7A}       % слот 3 графиков
\definecolor{cPanel}{HTML}{F1F5F9}    % заливка врезок
\definecolor{cGrid}{HTML}{CBD5E1}     % линии

\color{cInk}

% --- Заголовки ---
\titleformat{\section}{\Large\bfseries\color{cAccent}}{\thesection}{0.6em}{}
\titleformat{\subsection}{\large\bfseries\color{cAccent}}{\thesubsection}{0.6em}{}
\titlespacing*{\section}{0pt}{1.4\baselineskip}{0.5\baselineskip}

% --- Колонтитулы ---
\setlength{\headheight}{14pt}   % 12pt мало для 11pt-кегля: fancyhdr ругается
\pagestyle{fancy}
\fancyhf{}
\fancyhead[L]{\small\color{cInk2}\nouppercase{\leftmark}}
\fancyhead[R]{\small\color{cInk2}\thepage}
\renewcommand{\headrulewidth}{0.4pt}
\renewcommand{\headrule}{\color{cGrid}\hrule height 0.4pt}

% --- Подписи к рисункам и таблицам ---
\captionsetup{font=small,labelfont={bf,color=cAccent},margin=10pt}

% =====================================================================
%  Учебные врезки. Педагогический мандат CHARTER §11: интуиция ДО
%  формализма, предупреждения — выделены, фактура — со ссылкой на артефакт.
% =====================================================================

% Интуиция / пояснение «на пальцах»
\newtcolorbox{intuition}[1][]{
  enhanced, breakable, colback=cPanel, colframe=cS1,
  boxrule=0pt, leftrule=3pt, arc=1pt,
  fonttitle=\bfseries\color{cAccent},
  title={Интуиция#1}, coltitle=cAccent,
  attach boxed title to top left={xshift=6pt,yshift=-2pt},
  boxed title style={colback=cPanel,boxrule=0pt,arc=0pt},
  top=10pt
}

% Предупреждение / типичная ошибка
\newtcolorbox{warning}[1][]{
  enhanced, breakable, colback=cS2!6, colframe=cS2,
  boxrule=0pt, leftrule=3pt, arc=1pt,
  fonttitle=\bfseries, title={Осторожно#1}, coltitle=cS2,
  attach boxed title to top left={xshift=6pt,yshift=-2pt},
  boxed title style={colback=cS2!6,boxrule=0pt,arc=0pt},
  top=10pt
}

% Реальное измерение как иллюстрация тезиса. Учебных примеров «из головы» в курсе
% нет: каждое число здесь получено на настоящем приборе и настоящем образце.
\newtcolorbox{example}[1][]{
  enhanced, breakable, colback=cS3!6, colframe=cS3,
  boxrule=0pt, leftrule=3pt, arc=1pt,
  fonttitle=\bfseries, title={Из эксперимента#1}, coltitle=cS3!70!black,
  attach boxed title to top left={xshift=6pt,yshift=-2pt},
  boxed title style={colback=cS3!6,boxrule=0pt,arc=0pt},
  top=10pt
}

% Строгая формулировка после интуитивной
\newtcolorbox{strict}[1][]{
  enhanced, breakable, colback=white, colframe=cGrid,
  boxrule=0.8pt, arc=1pt,
  fonttitle=\bfseries\color{cInk}, title={Строго#1},
  attach boxed title to top left={xshift=6pt,yshift=-2pt},
  boxed title style={colback=white,boxrule=0pt,arc=0pt},
  top=10pt
}

% =====================================================================
%  Макросы обозначений. Единая нотация всего курса — менять только здесь.
% =====================================================================
\newcommand{\Deff}{D_{\mathrm{eff}}}
\newcommand{\Dphys}{D_{\mathrm{phys}}}
\newcommand{\tperp}{t_{\perp}}
\newcommand{\tpar}{t_{\parallel}}
\newcommand{\tleak}{\tau_{\mathrm{leak}}}
\newcommand{\dAIC}{\Delta\mathrm{AIC}}
\newcommand{\Efield}{\tilde{E}}
\newcommand{\jj}{\mathrm{j}}                      % инженерная мнимая единица
\newcommand{\Rot}[1]{R(#1)}
\DeclareMathOperator{\Real}{Re}
\DeclareMathOperator{\Imag}{Im}
\DeclareMathOperator{\cond}{cond}
\DeclareMathOperator*{\argmin}{arg\,min}
\DeclareMathOperator{\AIC}{AIC}

% Термин при первом употреблении
\newcommand{\term}[1]{\textbf{\color{cAccent}#1}}

% =====================================================================
%  Прослеживаемость числовых примеров.
%
%  В ОСНОВНОМ ТЕКСТЕ курса нет ни путей к файлам, ни имён программ, ни
%  внутренних обозначений наборов данных: измерение описывается физически
%  («плотная решётка D/P = 0.71»). Служебное приложение «Происхождение
%  числовых примеров» связывает каждый пример с породившим его файлом
%  результатов — это нужно для внутренней проверки, а не для читателя.
%
%  Публикация вовне: заменить \provenancetrue на \provenancefalse — приложение
%  и ссылки на него исчезнут, основной текст не изменится.
% =====================================================================
\newif\ifprovenance
\provenancetrue

% Метка источника — употребляется ТОЛЬКО внутри служебного приложения.
% \sloppy и мелкий кегль: пути к файлам длинные и в узкую колонку таблицы
% иначе не влезают (моноширинный шрифт не переносится по слогам).
\newcommand{\artifact}[1]{{\scriptsize\color{cInk2}\ttfamily\hbadness=10000 #1}}

% Ключ примера: якорь, по которому пример находится в таблице приложения
\newcommand{\srckey}[1]{\ifprovenance{\,\scriptsize\color{cInk2}[\,#1\,]}\fi}
